\section{Limitations and what we withdrew}
\label{sec:limitations}

\begin{enumerate}[leftmargin=*,itemsep=1pt]
\item \textbf{One trained setting.} E5 covers a single toy task at $d=\EfiveD$, $F=\EfiveF$,
with \EfiveSeeds{} seeds per condition. Nothing here is a claim about large trained language
models.
\item \textbf{Out-of-distribution evaluation.} The diagnostic is defined on Boolean sparse
states, while the toy models are trained on continuous sparse inputs. It therefore measures
which interface the learned code \emph{supports}, not what the network does on its native
distribution.
\item \textbf{Constant sparsity in the application.} The H\"{a}nni-template comparison assumes
$s=O(1)$, whereas the distributional results allow $s$ up to $O(d/\log d)$; real features are
continuous with heterogeneous sparsity.
\item \textbf{The floor bounds linear readouts only.} Nonlinear decoders and
distribution-specific average-case mechanisms are outside its scope --- which is the point of
the separation, but also a limit on what the diagnostic can certify.
\item \textbf{Fit with few points.} The $c\,d/\ln d$ fit in E6 rests on \EsixFitPoints{}
widths, on which a plain linear law fits nearly as well ($R^2=\EsixInterpRsqLin$ against
$\EsixInterpRsqLog$). The data are consistent with the predicted shape and do not discriminate
it from the alternatives; the logarithm's base is not identifiable at all.

\item \textbf{$s_{95}$ is a point estimate, not a certified threshold.} The trial budget shrinks
with width, and at $d=\EsixDmax$ even a perfect run gives a Wilson lower bound of
\EsixTopBestWilsonLow --- below the $95\%$ level being estimated. The reported thresholds are
best estimates; certifying them at confidence would need more trials at the large widths.

\item \textbf{The linear side's non-exactness is the theorem, not a measurement.} Wherever we
report that a linear readout ``cannot be exact'', that follows from
Theorem~\ref{thm:floor} by construction for any rank-$d$ unit-diagonal interface. What the
experiments measure is the \emph{magnitude} of that error for codes that were actually learned,
and how far a threshold decoder gets before it bites.
\item \textbf{Correlated curve points.} The nested-support trick makes each point of
$p_{\mathrm{rec}}(s)$ marginally unbiased but correlates points within a trial; the Wilson
intervals are per-point and should not be read as a joint confidence band.
\item \textbf{Open log gap.} The Adler--Shavit capacity bracket retains a single logarithmic
factor of slack, which is not resolved here.

\item \textbf{The leverage prediction is a sufficient condition, and its selection heuristic decays
with width.} Corollary~\ref{cor:failthresh} predicts the $L^2$ collapse from leverage alone and
holds at all three widths of Section~\ref{sec:res-primary}. But it is applied at $h_{\min}$, and it
is sufficient rather than necessary, which shows up in where the true worst feature sits: the
lowest-leverage window contains the true $\arg\min\kappa$ in \FullLfourTwoHundredArgminFound{} of 20
$L^4$ models at $d=200$ and \SubsetFourHundredFound{} of \ScSeeds{} at $d=400$, where the true
$\arg\min$ sits at leverage rank \SubsetFourHundredRank{} by median --- against every model of every
untrained code, at every width. Training pulls the two apart, and the gap widens with width, so the
theorem localises the failure less and less well as the code improves. Nothing reported depends on
the heuristic --- the frontiers quoted are exact --- but a reader should not take low leverage as a
reliable pointer to the binding feature, and past $d=200$ should not use it as a sampling window at
all.

\item \textbf{The exact frontier does not extrapolate much further.} Every $\kappa_i$ reported here is
exact --- all $F$ features of every model at all four widths, including the \ScSeeds{} models per arm
at $d=400$, $F=800$. An earlier version of this limitation said the $d=400$ sweep would be
unaffordable and that a fourth width would have to return to a subset estimate; it was run instead,
in a few hours on ten CPU workers, and this entry is corrected rather than removed because the
estimate was wrong in the direction that mattered. What remains true is the shape of the cost: one
linear programme per feature --- \FrontierLpSecondsTwoHundred\,s at $d=200$, from the recorded sweep
times --- growing steeply in $d$ while $F$ grows with it, with the wall-clock of every campaign in
Appendix~\ref{app:cost}. A fifth width would be a different
proposition, and the honest statement is that the statistic is exact everywhere we report it and that
we do not know where it stops being affordable.

\item \textbf{The nonlinear witness is built but not yet used.} The extractor of
Section~\ref{sec:witness} is implemented and validated --- every witness it returns is confirmed
inseparable by a linear programme, and on small codes against brute-force enumeration --- but two
things are missing. The Radon compression to $d+2$ states is not implemented, so the witnesses are
valid and not minimal; and no result reported in this paper is derived from one, so the sharper
claim it licenses is not made here.

\item \textbf{How much the ReLU costs is measurable in only three cells.} The decoder comparison
says the ReLU costs an $L^4$ code almost nothing and an untrained code a great deal
(Section~\ref{sec:res-primary}). Quantifying that on the frontier rather than on a fitted probe is
harder: $\kappa$ is defined on $\Phi b$, and $\mathrm{ReLU}(\Phi b)$ is not a linear image of the
state polytope, so the post-ReLU side must be sampled and tested rather than solved. Sampling can
only refute separability, never confirm it, so it overestimates the frontier --- here by
\ReluLooseMin{} to \ReluLooseMax{} levels against the exact $\kappa$ --- and the overestimate cancels
only in a difference taken with the same sampler and the same draws on both sides
(\ReluGapDraws{} draws per side, \ReluGapFeatures{} features per cell, sampler ceiling
$s\le\ReluGapSmax$).

Three of the nine cells give a measurable difference and they agree with the decoder comparison: at
$d=50$ the ReLU costs the $L^4$ code \ReluLfourFiftyGap{} levels
$[\ReluLfourFiftyCiLow,\ReluLfourFiftyCiHigh]$ on \ReluLfourFiftyPositive{} of \ReluGapFeatures{}
features and the untrained code \ReluRandFiftyGap{} $[\ReluRandFiftyCiLow,\ReluRandFiftyCiHigh]$ on
\ReluRandFiftyPositive{}; at $d=100$ the untrained code loses \ReluRandHundredGap{}
$[\ReluRandHundredCiLow,\ReluRandHundredCiHigh]$ on all \ReluRandHundredPositive{}. In
\ReluCensoredCells{} cells both sides reach the sampler ceiling, so the difference there is censoring
and not measurement, and we report no value --- including for every $L^4$ cell above $d=50$, which is
precisely where we would most want one. Every $L^2$ cell is a flat zero because its frontier is at
$1$ on both sides. A larger sampling budget would lift the ceiling and we have not spent one, so the
claim that the cost is small on trained codes rests on $d=50$ plus the decoder comparison, and not on
the frontier at the two larger widths.

\item \textbf{One control the analysis wants and does not have.} A random \emph{tight frame}, whose
$R_{\mathrm{geom}}$ is exactly $1$ rather than the \SaeCtlRgeomMax{} an i.i.d.\ code reaches at
worst, would sharpen
the claim that near-floor geometry is generic: as it stands the comparison is against a code that
is close to the floor rather than on it, so a small residual advantage for the trained code could be
conditioning rather than structure.

\item \textbf{The margin over an untrained code stops growing inside our own range.} This is the
sharpest limitation on the headline claim, and it comes from our own data rather than from something
we did not run. Section~\ref{sec:res-primary} reports the trained frontier at $d^{\ScExpTrained}$ and
the untrained one at $d^{\ScExpRandom}$: the untrained exponent is larger, the ratio falls at every
step, and the margin in sparsity levels peaks at $d=200$ and falls by \MarginDrop{}
$[\MarginDropCiLow,\MarginDropCiHigh]$ at $d=400$. ``Training buys the code'' therefore holds at every
width we measured and is not established as a property that survives scale. Two candidate
explanations are not separated here: a real ceiling on what the frontier can buy, or the fixed
optimisation budget of \EsevenSteps{} steps at every width, which gives the largest models the least
optimisation per parameter. Distinguishing them needs a budget sweep at fixed width, which is cheap,
and a fifth width, which is not.
\end{enumerate}

